% !TEX root = ../../main.tex
\chapter{Vorfeldanalyse (\Luca)}
\label{chap:vorfeldanalyse}

% ---------------------------------------------------------------
\section{Anforderungen und Abgrenzung}


\subsection{Anforderungen an das Backend}
Das Backend spielt in diesem Projekt eine zentrale Rolle. Seine Hauptaufgabe besteht in der Abbildung der Kernlogik und der Verwaltung der Daten.
Die Anforderungen ergeben sich aus der Zielsetzung der Webanwendung.
Dabei wird zwischen funktionalen und nicht-funktionalen Anforderungen unterschieden.
Maximilian Lutz zitiert in seiner Arbeit \cite{lutz2022} aus dem Buch \cite{broy2021} die Definition von Anforderungen.
Funktionale Anforderungen beschreiben, wie das System genutzt werden kann,
und stellen die Funktionen an der Systemgrenze in den Mittelpunkt.
Nicht"=funktionale Anforderungen legen hingegen fest, wie das System seine
Leistung erbringt \parencite[S.~16]{lutz2022}.\\
\\
\textbf{Konkrete funktionale Anforderungen des Projekts sind:}
\par
\begin{itemize}
    \item Verwaltung von Benutzern, Lerninhalten, Lernkarten, Tests und Ergebnissen;
    \item Abbildung der Lernergebnisse und des Wiederholungsstatus jedes einzelnen Benutzers;
    \item Durchführung der Wiederholungsplanung und Wissensstandschätzung durch den adaptiven Lernalgorithmus auf Grundlage der gespeicherten Lernergebnisse;
    \item Bereitstellung einer strukturierten API für die Kommunikation mit dem Frontend und der Infrastruktur.
\end{itemize}
\textbf{Konkrete nicht"=funktionale Anforderungen des Projekts sind:}
\par
\begin{itemize}
    \item Skalierbarkeit: Das Backend muss in der Lage sein, eine wachsende Anzahl an Benutzern und Daten zu bewältigen;
    \item Modularität: Die Backend-Logik sollte in klar abgegrenzte Module unterteilt sein, um Wartbarkeit und Erweiterbarkeit zu gewährleisten;
    \item Sicherheit: Das Backend muss sicherstellen, dass Benutzerdaten geschützt sind und nur autorisierte Zugriffe erfolgen;
    \item Performance: Das Backend sollte eine kurze Reaktionszeit bieten, um eine gute Benutzererfahrung zu gewährleisten.
\end{itemize}

\subsection{Abgrenzung zu Frontend und Infrastruktur}
\Luca ist für die reine Entwicklung eines skalierbaren Backends zuständig.
``Skalierbar'' bedeutet in diesem Zusammenhang, dass das Backend modular aufgebaut ist und sich um neue Funktionen erweitern lässt, ohne bestehende Komponenten umzubauen.
Die Kommunikation mit dem Frontend wird mit \Jakob konzeptioniert und erfolgt über eine klar definierte API.
Die Kommunikation mit der Infrastruktur, wie z. B. Datenbanken oder Cloud-Diensten, wird mit \Jai konzeptioniert und über das ORM-Framework Django realisiert.
Die Datenverwaltung und die Abbildung der Kernlogik und Speicherung der Lernergebnisse sind die Hauptaufgaben des Backends.
Die Entwicklung eines adaptiven Lernalgorithmus ist eine weitere wichtige Aufgabe des Backends.
Die Twelve-Factor-Methodik betrachtet Dienste wie eine Datenbank als
angebundene Ressourcen, die über eine in der Konfiguration hinterlegte Adresse
und Zugangsdaten angesprochen werden. Eine lokale Datenbank lässt sich so
durch eine extern verwaltete ersetzen, ohne den Anwendungscode zu ändern
\parencite[vgl.][]{wiggins12factor}.
Die Trennung der Zuständigkeiten von Benutzeroberfläche und Datenhaltung
macht die Oberfläche auf mehreren Plattformen besser einsetzbar und erleichtert
die Skalierung, weil die Serverkomponenten einfacher werden
\parencite[vgl.][Abschn.~5.1.2]{fielding2000}.
\Jakob ist zuständig für die Entwicklung der Benutzeroberfläche sowie Umsetzung der Benutzerinteraktion für Lernkarten, Tests und die Darstellung des individuellen Lernfortschritts.
\Jai ist zuständig für die Konzeption und Umsetzung einer skalierbaren Serverinfrastruktur und Datenbankinfrastruktur mit Active Directory, Lerninhalte-Import/-Export sowie Durchführung von Sicherheits- und Performancetests.

%todo:zitate sinnvoll einfügen, abgrenzen, was jai und jakob machen.

% ---------------------------------------------------------------
\section{Softwarearchitektur und Backend-Technologien}
\label{sec:architektur}
% - Monolith, modularer Monolith oder Microservices? Belege überprüfen!
% - Warum Django, welche Alternativen (z. B. FastAPI, Flask)?
% - Was bedeutet "skalierbar" und "modular" konkret hier?

% TODO: Belege überprüfen!

\subsection{Architekturansätze im Vergleich}
Relevante Architekturansätze für die Entwicklung des Backends sind der Monolith, der modulare Monolith und Microservices.
Die Vor- und Nachteile der drei Ansätze werden in Tabelle~\ref{tab:architekturvergleich} gegenübergestellt.\\ \\
\begingroup
\small
\setlength{\tabcolsep}{2pt}
\begin{longtable}{
  >{\raggedright\arraybackslash}p{2.4cm}
  >{\raggedright\arraybackslash}p{2.7cm}
  >{\raggedright\arraybackslash}p{2.9cm}
  >{\raggedright\arraybackslash}p{3.1cm}
  >{\raggedright\arraybackslash}p{1.8cm}}

\caption{Vergleich von Monolith, modularem Monolith und Microservices bei einer Neuentwicklung}
\label{tab:architekturvergleich} \\

\toprule
\textbf{Kriterium} & \textbf{Monolith} & \textbf{Modularer Monolith} & \textbf{Microservices} & \textbf{Beleg} \\
\midrule
\endfirsthead

\multicolumn{5}{l}{\small\itshape Fortsetzung von Tabelle~\ref{tab:architekturvergleich}} \\
\toprule
\textbf{Kriterium} & \textbf{Monolith} & \textbf{Modularer Monolith} & \textbf{Microservices} & \textbf{Beleg} \\
\midrule
\endhead

\midrule
\multicolumn{5}{r}{\small\itshape weiter auf der nächsten Seite} \\
\endfoot

\bottomrule
\endlastfoot

Deployment
  & Eine Einheit
  & Eine Einheit
  & Pro Service einzeln deploybar
  & \cite{zenodo2025petclinic}, \cite{faustino2022} \\

Code-Struktur
  & Eine Codebasis, Zugriffe überall möglich
  & Eine Codebasis, Module mit Schnittstellen
  & Getrennte Codebasen bzw. Projekte
  & \cite{faustino2022}, Eig. \\

Kopplung
  & Oft eng (direkter Zugriff auf alle Modelle)
  & Lose, Zugriff nur über Schnittstellen und DTOs
  & Lose, nur über API oder Events
  & \cite{faustino2022} \\

Datenhaltung
  & Eine Datenbank
  & Eine Datenbank, aber getrennte Teilmodelle
  & Meist je Service eine eigene Datenbank
  & \cite{faustino2022} \\

Kommunikation
  & Funktionsaufrufe im Prozess
  & Funktionsaufrufe im Prozess (über Schnittstelle)
  & HTTP/REST oder Message Broker
  & \cite{faustino2022} \\

Transaktionen
  & Normale ACID-Transaktionen
  & ACID möglich, da eine Datenbank
  & Keine verteilten ACID-Transaktionen, Eventual Consistency
  & \cite{faustino2022} \\

Latenz
  & Am niedrigsten
  & Nahe am Monolithen, abhängig von der DTO-Menge
  & Am höchsten (Netzwerk, Serialisierung)
  & \cite{faustino2022}, \cite{zenodo2025petclinic} \\

Ressourcenbedarf
  & Gering
  & Gering
  & Höher (mehr Prozesse, mehr Speicher)
  & \cite{zenodo2025petclinic} \\

Konfiguration und Betrieb
  & Einfach
  & Einfach
  & Aufwendig (mehrere Dienste und Konfigurationsdateien)
  & \cite{zenodo2025petclinic} \\

Anfänglicher Entwicklungsaufwand
  & Am geringsten
  & Mittel (Schnittstellen und DTOs definieren)
  & Am höchsten (APIs, Broker, Konsistenz, Deployment)
  & Eig., \cite{faustino2022} \\

Modulgrenzen später ändern
  & Leicht, aber ohne erzwungene Grenzen
  & Mittel (Refactoring im selben Prozess)
  & Schwer (API-Änderungen über Prozessgrenzen)
  & \cite{faustino2022}, Eig. \\

Gefahr falscher Modulgrenzen
  & Gering, da keine festen Grenzen
  & Mittel
  & Hoch
  & \cite{faustino2022} \\

Team-Unabhängigkeit
  & Gering
  & Mittel bis gut
  & Am besten (eigene Services pro Team)
  & \cite{faustino2022} \\

Skalierung
  & Gesamte Anwendung als Ganzes (z.\,B. mehr Worker)
  & Wie Monolith
  & Einzelne Services unabhängig skalierbar
  & Eig. \\

Ausfallsicherheit
  & Ein Fehler kann die ganze Anwendung betreffen
  & Wie Monolith
  & Ausfall eines Services betrifft nicht automatisch alle
  & \cite{aldebagy2018} \\

Testbarkeit
  & Einfach (ein Prozess)
  & Einfach, Module einzeln testbar
  & Aufwendiger (Integrationstests über Services)
  & Eig. \\

Typisches Risiko
  & \enquote{Big Ball of Mud}, wenn keine Struktur eingehalten wird
  & Grenzen werden im Alltag aufgeweicht
  & Overengineering, Betriebsaufwand
  & Eig. \\

\end{longtable}
\endgroup

\noindent\textit{Anmerkung:} \enquote{Eig.} kennzeichnet eigene Einschätzungen
ohne gesonderten Quellenbeleg.
\subsection{Auswahl des Frameworks}
Relevante Frameworks für die Entwicklung des Backends sind Django, FastAPI und Flask.
Die drei Frameworks werden miteinander verglichen. Die Vor- und Nachteile der drei Frameworks werden in Tabelle~\ref{tab:frameworkvergleich} gegenübergestellt.\\ \\


\begin{footnotesize}
\setlength{\tabcolsep}{3.5pt}
\renewcommand{\arraystretch}{1.15}

\begin{longtable}{
  >{\raggedright\arraybackslash}p{2.1cm}
  >{\raggedright\arraybackslash}p{3.4cm}
  >{\raggedright\arraybackslash}p{3.4cm}
  >{\raggedright\arraybackslash}p{3.4cm}
  >{\raggedright\arraybackslash}p{2.3cm}}

\caption{Vergleich von Django (mit Django REST Framework), FastAPI und Flask für das Backend}
\label{tab:frameworkvergleich} \\

\toprule
\textbf{Kriterium} & \textbf{Django (mit DRF)} & \textbf{FastAPI} & \textbf{Flask} & \textbf{Beleg} \\
\midrule
\endfirsthead

\multicolumn{5}{l}{\itshape Fortsetzung von Tabelle~\ref{tab:frameworkvergleich}} \\
\toprule
\textbf{Kriterium} & \textbf{Django (mit DRF)} & \textbf{FastAPI} & \textbf{Flask} & \textbf{Beleg} \\
\midrule
\endhead

\midrule
\multicolumn{5}{r}{\itshape weiter auf der nächsten Seite} \\
\endfoot

\bottomrule
\endlastfoot

Charakter
  & Umfangreiches Framework; mitgeliefert u.\,a. ORM, Migrationen, Admin-Oberfläche, Templates, Caching. API-Funktionen über das Django REST Framework (DRF)
  & Baut auf Starlette und Pydantic auf; will gezielt ein \enquote{Micro-Framework} sein
  & \enquote{Microframework}: schlanker, erweiterbarer Kern, trifft wenige Vorgaben (z.\,B. zur Datenbank)
  & \cite{djangoOverview}, \cite{drfHome}; \cite{fastapiAlternatives}; \cite{flaskForeword} \\

Datenbankzugriff
  & Eigener ORM mit Migrationen; laut FastAPI-Doku eng an relationale Datenbanken gekoppelt
  & Kein eigener ORM; ORMs und ODMs von Drittanbietern auf Pydantic-Basis nutzbar
  & Keine Datenbankintegration im Kern; Anbindung über Erweiterungen (z.\,B. Flask-SQLAlchemy)
  & \cite{djangoOverview}, \cite{fastapiAlternatives}; \cite{fastapiFeatures}; \cite{fastapiAlternatives}, \cite{bednarz2025} \\

Admin-Oberfläche
  & Wird aus den registrierten Modellen automatisch erzeugt
  & Nicht Teil des Kerns
  & Nicht Teil des Kerns
  & \cite{djangoOverview}; Eig. \\

Authentifizierung und Sicherheit
  & Benutzer- und Rechteverwaltung über \texttt{django.contrib.auth}; DRF ergänzt Authentifizierung, Berechtigungen und Throttling. Eingebauter Schutz u.\,a. gegen XSS, CSRF, SQL-Injection und Clickjacking
  & Integrierte Werkzeuge für HTTP Basic, OAuth2 (mit JWT) und API-Keys
  & Nicht im Kern; per Plug-ins ergänzbar
  & \cite{drfHome}, \cite{djangoSecurity}; \cite{fastapiFeatures}; \cite{fastapiAlternatives} \\

API-Dokumentation und Validierung
  & Serializer für ORM- und Nicht-ORM-Daten, \enquote{Browsable API}, Schema-Erzeugung als DRF-Funktion
  & Automatische interaktive Dokumentation (Swagger UI, ReDoc) auf Basis von OpenAPI; Validierung über Pydantic und Typ-Annotationen
  & Validierung und OpenAPI-Dokumentation über Zusatzpakete (z.\,B. Flask-apispec, Marshmallow, Webargs)
  & \cite{drfHome}; \cite{fastapiFeatures}; \cite{fastapiAlternatives} \\

Asynchronität
  & Async-Views; ORM, Cache, Authentifizierung und Sessions mit Async-Schnittstellen. Volle Wirkung unter ASGI, unter WSGI mit Zusatzkosten. Async-Unterstützung von DRF nicht geprüft
  & Auf ASGI (Starlette, Uvicorn) ausgelegt, nativ asynchron
  & Async-Views seit Version 2.0, aber unter WSGI bindet jeder Request einen Worker; für überwiegend asynchronen Code verweist die Doku auf Quart
  & \cite{djangoAsync}; \cite{fastapiAlternatives}, \cite{bednarz2025}; \cite{flaskAsync} \\

Durchsatz (Benchmark)
  & GET: 1\,660\,req/s; Echo: 1\,551\,req/s
  & GET: 4\,783 (SQLite) bzw. 4\,417 (SQLModel)\,req/s; Echo: 4\,024 bzw. 3\,627\,req/s
  & GET: ca.\ 2\,700\,req/s; Echo: 2\,194 (SQLite) bzw. 2\,143 (SQLAlchemy)\,req/s
  & \cite{bednarz2025} \\

Antwortzeit (Benchmark)
  & Einzelner Datensatz: 13,4\,ms; alle Datensätze: 75,0\,ms; Anlegen: 99,6\,ms
  & Einzeln: 2,5 (SQLite) bzw. 7,6 (SQLModel)\,ms; alle: 67,4 bzw. 168,8\,ms; Anlegen: 78,6 bzw. 92,4\,ms
  & Einzeln: 5,0 (SQLite) bzw. 8,5 (SQLAlchemy)\,ms; alle: 28,2 bzw. 69,2\,ms; Anlegen: 111,3 bzw. 116,0\,ms
  & \cite{bednarz2025} \\

Projektstruktur
  & Anwendung in Apps gegliedert (Beispiel-App \enquote{news} in der Doku)
  & Mehrere Dateien über Router möglich
  & Keine feste Projektstruktur vorgegeben
  & \cite{djangoOverview}; Eig.; \cite{bednarz2025} \\

Genannter Haupteinsatz
  & Ursprünglich zur HTML-Erzeugung im Backend gedacht; mit DRF auch für APIs; geeignet für größere Projekte mit Komplettumfang
  & Moderne, performante APIs und Microservices
  & Kleine bis mittlere Projekte, Prototypen, volle Kontrolle
  & \cite{fastapiAlternatives}, \cite{djangoOverview}; \cite{bednarz2025} \\

\end{longtable}

\noindent\textit{Anmerkung:} Die Benchmark-Werte stammen aus \cite{bednarz2025}
(Windows\,10, SQLite, Gunicorn mit 2 Workern und 2 Threads, nicht auf das jeweilige
Framework optimiert; Django\,5.2.1 mit DRF, Flask\,3.1.0, FastAPI\,0.115.12).
\enquote{Eig.} kennzeichnet eigene Einschätzungen ohne gesonderten Quellenbeleg.
\end{footnotesize}

\subsection{Skalierbarkeit und Modularität}

% ---------------------------------------------------------------
\section{Backend-Logik und Datenverwaltung}
\label{sec:backend-logik}
% - Entitäten: Benutzer, Lerninhalt, Lernkarte, Test, Ergebnis
% - Abbildung von Lernergebnissen und Wiederholungsstatus pro Benutzer
% - Strukturierung in Django-Apps

% TODO: Entitäten und Beziehungen festlegen

\subsection{Fachliche Entitäten}
\subsection{Speicherung von Lernergebnissen und Wiederholungsstatus}
\subsection{Projektstruktur in Django-Apps}

% ---------------------------------------------------------------
\section{Adaptiver Lernalgorithmus}
\label{sec:algorithmus}
% Kernthema, hier am meisten Recherche einplanen.
% - Was bedeutet "adaptiv", was wird individuell angepasst?
% - Verfahren zur Wiederholungsplanung (z. B. Leitner-System, SM-2, FSRS)
% - Verfahren zur Wissensstandschätzung (z. B. Elo-ähnliche Modelle, IRT)
% - Lernforschung zum verteilten Wiederholen: nur mit belastbaren Quellen belegen
% - Realistischer Umfang bei 180 h, Prüfung der Wirksamkeit

% TODO: Quellen zur Lernforschung suchen und in literatur.bib eintragen

\subsection{Begriff und Ziel der Adaptivität}
\subsection{Grundlagen aus der Lernforschung}
\subsection{Verfahren zur Wiederholungsplanung}
\subsection{Verfahren zur Wissensstandschätzung}
\subsection{Vergleich und Auswahl}


\subsection{Bewertung der Wirksamkeit}

% ---------------------------------------------------------------
\section{API-Konzeption}
\label{sec:api}
% - REST oder GraphQL, Django REST Framework oder Alternativen?
% - Versionierung, Fehlerformat, Dokumentation (z. B. OpenAPI)
% - Abstimmung der Endpunkte mit dem Frontend und der Infrastruktur

% TODO: API-Stile vergleichen

\subsection{API-Stile im Vergleich}
\subsection{Gestaltung der Schnittstelle}
\subsection{Abstimmung mit Frontend und Infrastruktur}

% ---------------------------------------------------------------
\section{Sicherheit}
\label{sec:sicherheit}
% - Authentifizierung: Sessions, Token, JWT;
% - Rollen und Berechtigungen (Schüler, Lehrer, Admin)
% - Risiken laut OWASP Top 10 für ein Django-Backend
% - DSGVO bei Lernfortschrittsdaten von Schülern

% TODO: Anforderungen zur Authentifizierung mit Kumar abstimmen

\subsection{Authentifizierung und Autorisierung}

\subsection{Datenschutz}

% ---------------------------------------------------------------
\section{Teststrategie}
\label{sec:tests}
% - Unit-, Integrations- und API-Tests, Werkzeuge (z. B. Django-Tests, pytest)
% - Bewertung der Qualität des Lernalgorithmus

% TODO: Testarten und Werkzeuge festlegen

\subsection{Testarten und Werkzeuge}
\subsection{Testen des Lernalgorithmus}

% ---------------------------------------------------------------
\section{Fazit der Vorfeldanalyse}
\label{sec:fazit-vorfeld}
% - Abgeleitete Entscheidungen (Technologie, Architektur, Algorithmus)
% - Verbleibende Risiken und offene Fragen

% TODO: Erst nach Abschluss der Recherche schreiben

\subsection{Getroffene Entscheidungen}
\subsection{Risiken und offene Fragen}